\cotefile{10}
\scanpage{60}
% block: heading
\uline{Catégorie tannakienne} (sur un corps $k$) \uline{de} \deleted{\unsure{grou|pour}} \uline{lien un groupe diagonalisable} $G=D(\handwritten{\unsure{\Gamma|T}})$ \handwritten{$(\Leftrightarrow\mathcal{M}$} semi-simple, et pour tout objet simple (ou isotypique) $S$ de $\handwritten{\mathcal{M}}$, de rang $d(S)$, $\operatorname{End}(S)$ est de rang $d(S)^{\handwritten{2}}\handwritten{)}$.

% block: classification and first display
Soit $\handwritten{\mathcal{M}}$ cette c atégorie. Elle est définie à équivalence près (respectant structures et contraintes) par un élément\inserted{\handwritten{$\xi$}} de $H^{\handwritten{2}}(k,G)$. On a d'autre part
\[
H^{\handwritten{2}}(k,G)\handwritten{\simeq}\operatorname{Hom}(\handwritten{\Gamma},\operatorname{Br}(k))\ ,
\]
donc $\handwritten{\mathcal{M}}$ est définie essentiellement par l'homomorphisme \deleted{\unsure{cano|canon}}
\[
\handwritten{\deleted{\unsure{\xi}}u:\ \Gamma\longrightarrow}\operatorname{Br}(k)\ .
\]

% block: grading and characterization
\deleted{Pour tout} La catégorie $\handwritten{\mathcal{M}}$ est graduée de type $\handwritten{\Gamma}$, les objets homogènes de type $\handwritten{\Gamma}$ sont\inserted{\handwritten{les objets}} isotypiques, donc pour tout $i\handwritten{\in\Gamma}$, il y a exactement un objet simple, à isomorphisme près, $\unsure{\mathbb{E}|E}_i$, \deleted{de} homogène de degré $i$. (Il semble que ce soit une deuxième caractérisation des catégories tannakiennes de lien $D(\handwritten{\Gamma})$,\inserted{\handwritten{$\Leftrightarrow$ le \unsure{pr.} \unsure{direct|tensl|tensoriel} de deux \unsure{isotypiques} est isotypique}}) Le degré de cet élément est égal à l'ordre\inserted{$d(i)$} de $u(i)$ dans $\operatorname{Br}(k)$. Donc l'anneau $K(\handwritten{\mathcal{M}})$ est donné par
\[
\begin{aligned}
K(\handwritten{\mathcal{M}})\handwritten{\subset}\underline{Z}[\handwritten{\Gamma}]\ ,\\
\handwritten{\sum}n_i e^i\handwritten{\in}K(\handwritten{\mathcal{M}})\quad\handwritten{\Longleftrightarrow}\quad n_i\handwritten{\equiv}0\ (d(i))\quad\text{pour tout }i\handwritten{\in\Gamma}
\end{aligned}
\]
La base canonique de $K(\handwritten{\mathcal{M}})$ formée par les objets simples étant formée des $d(i)e^i$. Si $M_i$ est l'objet simple de \deleted{poids} degré $i$, alors $\operatorname{End}(M_i)$ est le corps gauche \deleted{de sur} de centre $k$ dont l'invariant \unsure{daas|dans} $\operatorname{Br}(k)$ est $u(i)$.

% block: cancelled paragraph
\deleted{En termes d'un système de générateurs $(\,)_i\,\unsure{i\ I|i\in I}$ de $I$, \deleted{et}(d'un système d'objets \deleted{simples}\inserted{isotypiques} $\unsure{M|\mathbb{M}}_i$ \deleted{\unsure{de type}} degré $i$)et des}

% block: free group case
Si $\Gamma=Z^{(I)}$,\inserted{i.e. $G=G_m$,} la donnée de $u$\inserted{i.e. de $\handwritten{\xi}$} revient à la donnée d'éléments $\handwritten{\xi}_i\handwritten{\in}\operatorname{Br}(k)$. Choisissant pour tout $i$ un objet isotypique non nul de degré $i$, soit $E_i$, et désignant par $A_i=\operatorname{End}(\handwritten{M}_i)$ l'algèbre d'Azumaya associée,\inserted{de classe $\handwritten{\xi}_i$,} on réconstitue (à équivalence unique près\deleted{)} respectant tout)
% archive sheet number: 59
\scanpage{61}
% block: reconstruction continued
$\handwritten{\mathcal{M}}$ en termes de la \unsure{famille|faamille} des $A_i$ \deleted{et des entiers $d_i=\operatorname{rang}E_i$} comme suit\deleted{e}. On considère d'abord la catégories sous-additive \deleted{engendrée} ayant comme objets les éléments $\handwritten{\underline{\unsure{\gamma|Y}}}$ de $\underline{N}^{(I)}$, ($\handwritten{\underline{\gamma}}=(\handwritten{\gamma}_i)_{i\handwritten{\in}I}$ étant interprété comme $\handwritten{\bigotimes}E_i^{\handwritten{\gamma}_i}$), \deleted{ou} avec $\operatorname{Hom}(\handwritten{\underline{\gamma}},\handwritten{\underline{\gamma}'})=0$ si $\handwritten{\underline{\gamma}\ne\underline{\gamma}'}$ et $=\bigotimes_{\handwritten{i}} A_i^{\otimes\handwritten{\gamma_i}}$ si $\handwritten{\underline{\gamma}=\underline{\gamma}'}$, la composition des morphismes évidentes. On définit de façon évidente le foncteur $\otimes$ (addition de $\underline{N}^{(I)}$ sur les objets, produit tensoriel sur les homomorphismes) et ses contraintes,\margin{\handwritten{\unsure{attention à} contrainte de commutativité}} de façon que l'on trouve une sous-$\otimes$-catégorie. Son enveloppe Karoubienne est une\inserted{$\otimes$-}catégorie abélienne, qui s'identifie à la sous-catégorie pleine de $\handwritten{\mathcal{M}}$ formée des éléments dont les composantes homogènes non nulles sont à degrés $\handwritten{\in}\underline{N}^{(I)}$,\inserted{et a un foncteur fibre sur \deleted{$k$}\inserted{$\overset{\handwritten{-}}{k}$.}} D'autre part, pour tout $i$, si $A_i$ est de \deleted{degré}\inserted{rang} $d_i^2$ (donc $E_i$ de rang $d_i$), soit $L_i=\det E_i={\handwritten{\unsure{A|\bigwedge}_{\unsure{2|k}}}}^{d_i}\handwritten{E_i}$ (où $E_i$ désigne l'objet de la catégorie construite correspondant \deleted{à \unsure{il|le}} $i$.ème élément de la base canonique de $\underline{Z}^{(I)}$, soit $1_i$ ...), c'est un objet préinversible ($\inserted{\handwritten{X\to}}\deleted{\unsure{\otimes}}L_i\otimes\handwritten{X}$ est pleinement fidèle) et en inversant les $L_i$, on trouve (à \deleted{\unsure{équ}}équivalence près\deleted{)} respectant tout) la catégorie $\handwritten{\mathcal{M}}$.

% block: nonfree case
Quand $\handwritten{\unsure{\Gamma|T}}$ n'est pas libre, on trouve une construction analogue en termes d'un système de générateurs\inserted{convenable} $(\handwritten{\gamma}_i)_{i\handwritten{\in}I}$, mais en``\unsure{renant|rendant} \deleted{de plus} égaux à 1 \deleted{x}'' de plus certains éléments.\margin{A expliciter !} Ces constructions sont bien adaptées pour décrire les $\otimes$-foncteurs de $\handwritten{\mathcal{M}}$ dans une \deleted{$\otimes$-}catégorie tannakienne, et pour décrire le changement de base.( Pour cette dernière raison notamment, il aurait été maladroit dans le cas général de prendre les $M_i$ simples ....). Attention, dans la description des $\otimes$-foncteurs en question, de tenir compte des contraintes, notamment de celle de commutativité; il semble qu'il doive revenir au même d'exiger que le rang est conservé ...
% archive sheet number: 60
\scanpage{62}
% block: title
Etude numérique des catégories tannakiennes.

% block: a absolutely semisimple objects
a) Soit $\handwritten{\mathcal{M}}$ une catégorie abélienne $k$-linéaire avec $\operatorname{Hom}$ de dimension finie ($k$ un corps). Si $M$ est un objet simple de $\handwritten{\mathcal{M}}$, la sous-catégorie abélienne de $\handwritten{\mathcal{M}}$ qu'il engendre est équivalente à la catégorie des vectoriels de dimension finie sur le corps $K=\operatorname{End}(M)$. On voit \unsure{ainss|ain\deleted{is}si|ain\deleted{si}is} que l'\unsure{\deleted{i}o|\deleted{o}i|\.{o}}bjet $M$ reste semi-simple par une extension $k'/k$ sss $K\otimes_k k'$ est un comp\unsure{\deleted{o}d|\deleted{d}o|d}sé \deleted{de corps} d'algèbres de matrices (i.e. est une algèbre semi-simple). Ceci est v\deleted{rai}\handwritten{\unsure{èrifié|érifié}} pour tout $k'$ sss c'est vrai pour \deleted{$k'$ clôture}\inserted{\handwritten{\uuline{une} ext.}} \uline{parfaite} de $k$, \deleted{\unsure{i|,}}ou sss le centre \deleted{de} $Z$ de $K$ est séparable sur $k$. On dit alors que $M$ est absolument semi-simple (à généraliser au cas n'est pas supposé simple).

% block: classes and base extension
Le groupe $K(\handwritten{\mathcal{M}})$ est le groupe libre engendré par l'ensemble $\handwritten{\Sigma}(\handwritten{\mathcal{M}})$ des classes d'objets simples de $\handwritten{\mathcal{M}}$. Si $k'$ est une extension de $k$, la fa\deleted{\unsure{y|c}}çon dont un objet simple $M$ de $\handwritten{\mathcal{M}}$ tel que $M_{k'}$ soit semi-simple se décompose enn objets simples se voit sur la structure de $K\otimes_k k'=K'$ : si $K'$ est le produit d'algèbres $M_{n_i}(K'_i)$, où les $K'_i$ sont des corps gauches,\inserted{$1\leq i\leq r$,} alors $M_{k'}$ se décompose en $r$ composantsisotypiques (correspondants aux $K'_i$) chacun ayant $n_i$ composants simples\inserted{$M'_i$}. Si $k'$ est une extension quasi-galoisienne de $k$, les classes des $M'_i$ sont conjuguées entre elles par l'action de $\operatorname{Gal}(k'/k)=\handwritten{\unsure{\pi|\Pi}}$. On trouve alors une bijection canonique (si l'hypothèse faite est vérifiée pour tout objet simple \deleted{de} $M$ de $\handwritten{\mathcal{M}}$, p.ex. si $k'/k$ est même galoisienne)
\[
\handwritten{\Sigma(\mathcal{M})\simeq\Sigma(\mathcal{M}')/\pi}\ .
\]
Il y aurait lieu de généraliser au cas où on ne suppose pas nécessairement que les $K\otimes_k k'$ sont semi-simples ....

% block: numerical data
\inserted{Si $\handwritten{\mathcal{M}}$ est semi-simple,} La catégorie $\handwritten{\mathcal{M}}$ est connue à équivalence près quand on connait l'ensemble $\handwritten{\Sigma}=\handwritten{\Sigma(\mathcal{M})}$, et l'application
\[
\handwritten{\sigma\longmapsto (Z_\sigma,\xi_\sigma\in}\operatorname{Br}(\handwritten{Z_\sigma})\handwritten{)}
\]
associant à toute classe\inserted{\handwritten{$\sigma$}} d'objet simple $\unsure{\mathbb{M}|M}$ de $S$ le \deleted{classe de}\inserted{\handwritten{centre $Z_\sigma$}} $\operatorname{End}(\handwritten{M})$
% archive sheet number: 61
\scanpage{63}
% block: numerical data continued
et la classe de $\operatorname{End}(\handwritten{M})$ dans $\operatorname{Br}(Z_{\handwritten{\sigma}})$ (le couple $(Z_{\handwritten{\sigma}},\handwritten{\xi_\sigma})$ étant défini à isomorphisme près). Sans supposer $\handwritten{\mathcal{M}}$ semi\handwritten{\unsure{(}}-simple, on peut\inserted{\handwritten{encore}} dire que la plus grande sous-catégorie semi-simple(engendrée par les objets semi-simples) est connue à équivalence près par la connaissance précédente, qui implique donc la connaissance de $K(\handwritten{\mathcal{M}})=\handwritten{\mathbb{Z}^{(\Sigma)}}$,\inserted{\handwritten{sa partie effective,}} et sa variance pour les changements de\deleted{s} base par $k'$ variable sur $k$. De façon précise,$\handwritten{\Sigma'=\Sigma(\mathcal{M}_{k'})}$ en tant qu'ensemble sur $\handwritten{\Sigma}$ correspond à l'ensemble \deleted{des corps} dont la fibre \deleted{est} en $\handwritten{\sigma}$ est l'ensemble des corps composants de $Z_{\handwritten{\sigma}}\otimes_k k'$, soient $Z'^{i}_{\handwritten{\sigma}}\deleted{\unsure{\times}}$, l'application $K(\handwritten{\mathcal{M}})=\underline{Z}^{(\handwritten{\Sigma})}\handwritten{\longrightarrow}K(\handwritten{\mathcal{M}_{k'}})=\underline{Z}^{(\handwritten{\Sigma'})}$ est donné par
\[
\handwritten{\sigma\longmapsto}\ n_i\ (\handwritten{\sigma},i)\ ,
\]
où $n_i$ est donné par la règle suivante: si $d$ est \deleted{le degré de} l'ordre de $\handwritten{\xi_\sigma\in}\operatorname{Br}(Z_{\handwritten{\sigma}})$, celui de son image dans $\operatorname{Br}(Z'^{i}_{\handwritten{\sigma}})$ est un \deleted{\unsure{idivs|divis}} diviseur $d'_i$ de $d$, et on a $n_i=d/d_i$. (A généraliser au cas où les $Z_{\handwritten{\sigma}}$ ne sont pas nécessairement séparables). Si $k'$ est une \deleted{\unsure{clô}} extension alg. close de $k$, l'ensemble $\handwritten{\Sigma'}$ s'identifie à l'ensemble somme des \uline{fibres} en $k'$ des $\operatorname{Spec}(Z_{\handwritten{\sigma}})$ ....

% block: b numerical characters
b) Supposons maintenant $\handwritten{\mathcal{M}}$ \deleted{semi-tannakienne} munie d'un produit tensoriel à contraintes d'associativité et commutativité, avec 1.

On désigne sous le nom de ``caractères numériques\deleted{''} de $\handwritten{\Sigma}$ '' les données suivantes:
\margin{\handwritten{préciser cette notion, en faisant d'un ``caractère numérique'' une structure susceptible de \unsure{variation} fonctorielle...}}

% block: numerical characters item 1
\unsure{à|1}) l'ensemble (à bijection près)\inserted{\handwritten{$\Sigma$}} des \deleted{\unsure{com}} classes d'isomorphie d'objets simples de $\handwritten{\mathcal{M}}$.

% block: numerical characters item 2
2) La donnée de la multiplication de $K(\handwritten{\mathcal{M}})={\unsure{\mathbb{Z}|B|Z}}^{(\handwritten{\Sigma})}$; ce qui revient à la donnée pour tout $\handwritten{\sigma,\tau\in\Sigma}$ du produit
\[
\handwritten{\sigma\tau}=\handwritten{\sum_\rho}c_{\handwritten{\sigma\tau}}^{\handwritten{\rho}}\handwritten{\rho}\qquad\handwritten{(c_{\sigma,\tau}^{\rho}\in\mathbb{N})},
\]
\handwritten{i.e. du système d'entiers naturels ${\unsure{\sigma|c}}_{\sigma,\tau}^{\rho}$.}
% archive sheet number: 62
\scanpage{64}
% block: numerical characters item 3
\handwritten{3}) La donnée de l'application
\[
\handwritten{\sigma\unsure{\mathrel{\vert\! -\! -}|\mapsto}}(Z_{\handwritten{\sigma}},\handwritten{\xi_\sigma})\ ,\ \handwritten{\xi_\sigma\in}\operatorname{Br}(Z_{\handwritten{\sigma}})
\]
envisagée dans a).

% block: numerical characters item 4
4) La donnée de la fonction
\[
\operatorname{rang}\ :\ \handwritten{\mathcal{M}\longrightarrow}\underline{Z}
\]
lorsque $\handwritten{\mathcal{M}}$ est supposée\deleted{semi-}tannakienne.\deleted{\unsure{(i.e.}}

% block: effective part
On notera que la donnée de 1),2) équivaut à la donnée de l'anneau commutatif $K=K(\handwritten{\mathcal{M}})$, avec sa ``partie effective
\[
\begin{aligned}
K^{\mathrm{eff}}&=\text{sous-groupe de }K(\handwritten{\mathcal{M}})\text{ engendré par les}\\
&\qquad c\handwritten{\ell}(M),\ M\handwritten{\in}\operatorname{Ob}\quad,\\
&=\underline{N}^{(\handwritten{\Sigma})}
\end{aligned}
\]
(NB $\handwritten{\Sigma}$ se récupère à partir de \deleted{\unsure{ces}}\inserted{\handwritten{la}} donnée\inserted{\handwritten{$K^{\mathrm{eff}}\subset K$}} comme l'ensemble des éléments non nuls minimaux de $K^{\mathrm{eff}}$). La donnée de 3) permet alors de trouver, comme explicité dans a), la variance fonctorielle de l'anneau $K(\handwritten{\mathcal{M}}_{k'})$ par rapport à $k'$. Enfin la donnée de 4) permet en principe de retrouver les polynômes caractéristiques (en particulier, la trace et le déterminant) d'un endomorphisme d'un objet \deleted{de}(semi-simple ?) $M$ de $E$, en termes de polynômes caractéristiques réduits (resp. trace réduite, norme réduite) dans les algèbres d'Azumaya\inserted{$\operatorname{End}(M_i)$,} correspondant aux composantes isotypiques $M_i$ de $M$. Pour $M$ isotypique de rang $n$, on aura en effet, si $K=\operatorname{End}(M)$ a comme centre $Z$ de rang $r$ sur $k$, donc $\unsure{m|n}=n'\unsure{\deleted{d}r|\deleted{r}d|r}$, et $K$ de rang $d^2$ sur $Z$, avec $d\handwritten{\mid}n'$ (NB $M$ de rang $n'$ sur $Z$):
\[
\begin{aligned}
\det_M f&=\deleted{\det}N_{Z/k}(\operatorname{Nr}_{K/Z}(f))^{n'/d}\\
\operatorname{Tr}_M f&=\operatorname{Tr}_{Z/k}(\operatorname{Trr}_{K/Z}(f))\ \frac{n'}{d}\\
\operatorname{Pol}_M(f,t)&=N_{Z/k}(\operatorname{Polr}_{K/Z}(f,t))^{n'/d}\quad\inserted{(n'/d=n/dr)}\ .
\end{aligned}
\]
\margin{\handwritten{\unsure{Signaler} que la connaissance des car. num. de $\mathcal{M}$ implique celle des car. numériques de $\mathcal{M}_{k'}$ pour \unsure{toute} extension $k'$ de $k$, (``fonctorielle en $k'$'')}}

% block: warning about reconstruction
On fera attention que la connaissance des caractères numériques de la catégorie tannakienne $\handwritten{\mathcal{M}}$ ne permet pas en général de réconstituer celle-ci à équivalence près, même si $k$ est un corps alg. clos de car. 0 (de sorte que la donnée de 3) devient vide) et $\handwritten{\mathcal{M}}$ semi-simple
% archive sheet number: 63
\scanpage{65}
% block: warning continued and exception
et même dans le cas particulier où\inserted{\handwritten{de plus}} $\handwritten{\mathcal{M}}$ est la catégorie des représentations d'un groupe fini (? ? à vérifier dans littérature qu'un groupe fini n'est pas déterminé par la connaissance de la table de multiplication de \deleted{\unsure{ses il|ses it}} son ensemble de caractères, et les rangs de ceux-ci ...) Comme exception à cette remarque, signalons cependant le cas où $\handwritten{\mathcal{M}}$ est tannakienne à lien diagonalisable (a\unsure{\deleted{n}l|n|l}ors $\handwritten{\Sigma}$ devient un sous-groupe de $K^{\unsure{\times|x|\mathbb{x}}}$, \deleted{donc}\inserted{qui} détermine le lien, et la donnée de 3) donne l'homomorphisme $\handwritten{\Sigma\longrightarrow}\operatorname{Br}(k)$ qui suffit à tout déterminer). Regarder le cas où le lien est un groupe de t.m.

% block: c multiplicative type
c) Cas d'un lien qui est de t.m. On suppose connu le caractère numérique \deleted{\unsure{(fonc}} de $\handwritten{\mathcal{M}}$. Prenant le changement de base par une clôture séparable $k'$ de $k$, on trouve que $\handwritten{\Sigma'=\Sigma(\mathcal{M}}_{k'}\deleted{\unsure{\times}}\handwritten{)}$ (canoniquement déduit de la donnée numérique, avec l'opération de $\handwritten{\pi}=\operatorname{Gal}(k'/k)$ dessus) est muni d'une structure de groupe induite par la multiplication de $K(\handwritten{\mathcal{M}}_{k'})=\underline{Z}^{(\handwritten{\Sigma'})}$, \deleted{\unsure{et au|et av}} respectée par l'action de $\handwritten{\pi}$. Ceci nous donne le lien $G=D(\handwritten{\Gamma})$ (où $\handwritten{\Gamma}$ est le groupe $\handwritten{\Sigma'}$ avec l'action de $\handwritten{\pi}$). L'ensemble $\handwritten{\Sigma}$ s'identifie à l'ensemble des orbites $\handwritten{\sigma}$ de $\handwritten{\pi}$ opérant sur $\handwritten{\Gamma}$. Si $\handwritten{\sigma}$ est une telle orbite, le corps $Z_{\handwritten{\sigma}}$ n'est autre que le corps des invariants $Z_{\handwritten{\sigma}}\handwritten{\subset}k'$ du stabilisateur d'un point de l'orbite $\handwritten{\sigma}$. Si $M_{\handwritten{\sigma}}$ est un objet de type $\handwritten{\sigma}$, on a un homomorphisme
\[
G\handwritten{\longrightarrow}\underline{\operatorname{Aut}}(M_{\handwritten{\sigma}})=\underline{K}_{\handwritten{\sigma}}^\times
\]
(les soulignés indiquent qu'on prend des schémas sur $k$ ...) qui se factorise par le centre $Z_{\handwritten{\sigma}}$ de $K_{\handwritten{\sigma}}$ en
\[
u_{\handwritten{\sigma}}:\ G\handwritten{\unsure{\mathrel{-\! -}|\longrightarrow}}\underline{Z}_{\handwritten{\sigma}}^\times\handwritten{=}{}_{Z}\handwritten{\prod}_{/k}\underline{G}_m\ ,
\]
et on vérifie aussitôt qu'on a
\[
u_{\handwritten{\sigma}}(\handwritten{\xi})=\handwritten{\xi_\sigma}\ ,
\]
où
% archive sheet number: 64
\scanpage{66}
% block: gerbe class
\[
\handwritten{\xi\in}H^{\handwritten{2}}(k_{\mathrm{pl}},G)
\]
est la classe de la gerbe tannakienne qui définit $\handwritten{\mathcal{M}}$. Donc la connaissance du caractère numérique de $\handwritten{\mathcal{M}}$\deleted{\inserted{(qui implique déjà la donnée}} équivaut à la donnée, en plus \unsure{du|de|da} groupe de t.m. $G$ sur $k$, des classes $u_{\handwritten{\unsure{*|\sigma}}}(\handwritten{\xi})$ pour tous les homomorphismes $u$ de $G$ dans des tores de la forme $\handwritten{\prod}_{Z/k}\underline{G}_m$, où $Z$ est une extension finie de $k$. \deleted{(}On vient de voir en effet que la donnée du caractère numérique permet de reconstituer ce qui précède. Pour voir l'inverse, notons que \deleted{si} le rang $\unsure{\deleted{d}n|\deleted{n}d|n|d}(\handwritten{\unsure{\sigma|\xi}})$ de $M_{\handwritten{\sigma}}$ est déterminé par la formule
\[
n(\handwritten{\sigma})=\operatorname{card}(\handwritten{\sigma})\operatorname{ord}(\handwritten{\xi_\sigma})
\]
(où $\operatorname{card}\handwritten{\sigma}$ est le cardinal de l'orbite $\handwritten{\sigma}$, $\operatorname{ord}(\handwritten{\xi_\sigma})$ est l'ordre de $\handwritten{\xi_\sigma}$ dans $\operatorname{Br}(Z_{\handwritten{\sigma}})$).

% block: determination question
La question de savoir si le caractère numérique détermine à équivalence de catégories tannakiennes près revient donc, pour $G$ fixé, à la question de savoir si un élément \deleted{de} $\handwritten{\xi\in}H^{\handwritten{2}}(k_{\mathrm{pl}},G_{\deleted{\times}})$ est connu quand on connait les $u_{\handwritten{*}}(\handwritten{\xi})$ pour $u:G\handwritten{\unsure{\mathrel{-\! -}|\longrightarrow}}\underline{Z}^{\unsure{\times|x|\mathbb{x}}}$ comme dessus. \deleted{\unsure{On va exes|On va ex}} C'est le cas pour $G$ diagonalisable \deleted{pro-dénombrable} (pro-dénombrable ?) On va examiner deux autres cas.

% block: example 1
\uline{Exemple} \unsure{$\mathbb{I}$|I|1}. $G$ est un tore de dimension 1, donc (s'il n'est trivial) tordu par une extension quadratique $Z$ de $k$. On a donc une suite exacte
\diagram{\begin{tikzcd}[ampersand replacement=\&]
0 \handwritten{\arrow[r]} \& G \handwritten{\arrow[r,dash]} \arrow[r,phantom,"u"] \& \underline{Z}^{\times} \handwritten{\arrow[r]} \arrow[r,phantom,"N_{Z/k}"] \& \underline{G}_m \handwritten{\arrow[r]} \& 0\ ,
\end{tikzcd}}
d'où on conclut, par la suite exacte de cohomologie et le th. 90 que $\operatorname{Ker}(u:H^{\handwritten{2}}(k,G)\handwritten{\longrightarrow}H^{\handwritten{2}}(k,\underline{Z}^{\times}))$ est nul. Donc dans le cas envisagé, le caractère numérique détermine $\handwritten{\mathcal{M}}$ à équivalence près.

% block: cancelled example 2
\deleted{\uline{Exemple} 2. Gardons \deleted{le} $G$ comme dessus, soit $n$ un entier $\quad2$, et soit $G'=\operatorname{Ker}n.\operatorname{id}_G$. La suite exacte}
\diagram{\begin{tikzcd}[ampersand replacement=\&]
\deleted{0} \& \deleted{G'} \& \deleted{G} \arrow[r,phantom,"\deleted{n}"] \& \deleted{G} \& \deleted{0}
\end{tikzcd}}
\deleted{et le résultat précédent montre que l'ensemble des classes de}
% archive sheet number: 65
\scanpage{67}
% block: title
\uline{Catégories tannakiennes définies par des cristaux.}

% block: crystals item 1
\handwritten{1.} Soit $k$ un corps de car. $p>0$, q\handwritten{u'}on regarde comme algèbre sur $\underline{Z}_p$, dont l'idéal maximal est muni de puissances divisées. Cela donne un sens au site cristallin de $k$ (sur $\underline{Z}_p$, q\handwritten{u}alifié aussi de ``absolu''), et aux Modules loc. libres de type fini\inserted{(resp. de p\unsure{\deleted{t}r|r|t}ésentation finie)} sur ledit, qu'on appellera aussi cristaux en modules (localement libres\inserted{resp. de prés; finie}) sur $k$. Ces cristaux forment une $\otimes$-catégorie $\underline{Z}_p$-linéaire, \deleted{avec des $\operatorname{Hom}$ qui sont des modules de type fini sur $\underline{Z}_p$} On peut expliciter cette catégorie à l'aide d'un $p$-anneau $W$ de corps résiduel $k$, comme la catégorie des modules libres de type fini (resp. de type fini) $M$ sur $\unsure{\mathbb{W}|W}$, avec pour tout entier $n$ une connexion absolue sur $M_n=M/p^nM$ (relativement à $W_n$ sur $\underline{Z}/p^n\underline{Z}$), à courbure nulle, ces connexions se recollant pour $\unsure{p|n}$ variable. Lorsque $\unsure{\deleted{p}k|k|p}$ est parfait, $W$ est l'anneau des vecteurs de Witt construit sur $k$, et dans la description précédente, on peut oublier les connexions: les catégories envisagées sont simplement la catégorie des modules libres de t.f (resp. des modules de t.f.) sur $W$. Si $K$ est le corps des fractions de $W$, la catégorie des \uline{isocristaux} \deleted{sur} (=cristaux à isogénie près) sur $k$ est équivalente à $\operatorname{Modf}(K)$. \deleted{C'est une catégorie tannakienne sur $\underline{Q}_p$}

% block: crystals item 2
\handwritten{2.} La catégorie des cristaux sur $k$ dépend fonctoriellement de $k$, d'où pour tout cristal la possibilité de définir son transformé par le changement de base $\handwritten{\sigma}_k:x\handwritten{\mapsto}x^p:\ k\handwritten{\longrightarrow}k$, soit $M^{\handwritten{\unsure{\sigma|\sigma'}}}$. La donnée d'un $M$ avec un hèmomorphisme
\[
F_M:\ M^{\handwritten{\sigma}}\handwritten{\longrightarrow}M
\]
s'appelle un $F$\handwritten{\unsure{(}}-\uline{cristal sur $k$}. On notera que les $F$-cristaux forment une\inserted{$\otimes$-}catégorie\inserted{$\underline{Z}_p$-linéaire} dont les $\operatorname{Hom}$ sont des modules de type fini sur $\underline{Z}_p$ (attention, pas\inserted{\handwritten{des modules}} sur $W$, même si $k$ est parfait !). Ce n'est pas une \deleted{\unsure{catégorie|ctatégorie}} catégorie tannakienne, même une fois tensorisée par $\underline{Q}_p$, car il manque
% archive sheet number: 66
\scanpage{68}
% block: nondegeneracy
un $\underline{\operatorname{Hom}}$ interne. On dit qu'un $F$-cristal $M$ est \uline{non dégénéré} s'il existe un entier $i\handwritten{\geq}0$ et un homomorphisme $V:M\handwritten{\longrightarrow}M^{\handwritten{\sigma}}$ tel que $VF=p^i$. Soit $\operatorname{FCriso}^{\inserted{\handwritten{+}}}(k)$ la catégorie \deleted{des isocristaux} déduite de la catégorie des \deleted{cr} $F$-cri\deleted{z}staux non dégénérés en localisant par rapport à $p$ ; c'est donc une sous-$\otimes$-catégorie de celle des $F$-cristaux à isogénie près. On l'appellera la catégorie des $F$-\uline{isocristaux effectifs}, et on laissera tomber le qualificatif ``iso'' s'il n'y a pas de risque de confusion.

% block: Tate crystal and inversion
\uline{Le} $F$-\uline{isocristal de Tate\inserted{\handwritten{\ }}inverse}\inserted{$K(-1)$} est déduit du $F$\handwritten{\unsure{(}}-cristal\inserted{$W(-1)$} défini par le cristal \uline{unité}, avec $F_{K(-1)}=p.\operatorname{id}$ (cet $F$-cristal est évidemment non dégénéré !). On voit alors que $K(-1)$ est faiblement inversible, et on constate que lorsqu'on ajoute formellement à $\mathrm{FC}\overset{\handwritten{\unsure{\cdot\ \cdot}}}{\underset{r}{r}}\mathrm{iso}^{+}(k)$ \deleted{des} un in\handwritten{v}erse $K(1)$ de $K\handwritten{(-1),}$ alors on trouve une catégorie tannakienne sur $\underline{Q}_p$, ayant un foncteur fibre évident sur $K$; cette catégorie est notée $\operatorname{FCriso}(\unsure{k|h})$, et appelée \uline{catégorie des} $F$-\uline{isocristaux} \handwritten{(}ou simplement, s'il n'y \deleted{pa} a pas risque de confusion, des $F$-cristaux) sur $k$. $K(1)$ est le $F$\handwritten{(}-cristal de Tate, et on désigne par $M\handwritten{\mapsto}M(n)$ le foncteur de tensorisation par sa puissance tensorielle $n$.ème. Si $n=-m$, $m\handwritten{\geq}0$, ce fonc\deleted{\unsure{teur}}teur est induit par le foncteur sur les $F$-cristaux effectifs (sans iso\unsure{)} consistant à multiplier $F$ par $p^m$.

% block: crystals item 3
\handwritten{3.} Lorsque $k$ est parfait, désignant par $\handwritten{\sigma}$ l'automorphisme de Frobenius de $W=W(k)$, la donnée d'un $F$-cristal revient à la donnée d'un module de type fini $M$ sur $W$ muni d'un homomorphisme
\[
F_M\handwritten{\unsure{\div|:|=}}\quad M^{\handwritten{\sigma}}=M\otimes_W(W,\handwritten{\sigma})\handwritten{\longrightarrow}M\ ,
\]
et le $F$\handwritten{(}-cristal est non dégénéré sss l'homomorphisme précédent est une isogénie i.e. est de rang égal à $\operatorname{rang}M$, ou encore sss il induit un isomorphisme après tensorisation par $K$. La catégorie \deleted{$\operatorname{Fiso}$} $\mathrm{Fc}\handwritten{ri}\mathrm{s}(k)$ s'identifie alors à la catégorie des vectoriels de dimension finie
% archive sheet number: 67
